\chapter{Происхождение родительского функционала разрывов и проводимостей KMS}
\label{ch:v9-endpoint-kms-parameter-origin}

\section{Шесть параметров после KMS completion}

Предыдущий гейт оставил две тройки положительных данных
\begin{equation}
 \boldsymbol\theta=(\theta_s,\theta_a,\theta_t),
 \quad \theta_\alpha=\beta\Delta_\alpha>0,
 \qquad
 \boldsymbol\kappa=(\kappa_s,\kappa_a,\kappa_t)>0,
 \label{eq:v9-kms-origin-six-parameters}
\end{equation}
причём $q_\alpha=e^{-\theta_\alpha}$. Вектор $\boldsymbol\theta$ задаёт
stationary Gibbs weights, а $\boldsymbol\kappa$ задаёт symmetric edge
conductances и темп приближения к stationarity.

Channel multiplicity, grading и charge дают три строки
\begin{equation}
 B_{\rm type}=
 \begin{pmatrix}
 1&1&3\\
 1&-1&3\\
 0&0&-3
 \end{pmatrix},
 \qquad \det B_{\rm type}=6,
 \qquad \rank B_{\rm type}=3.
 \label{eq:v9-kms-origin-type-matrix}
\end{equation}
Следовательно, типы различают $s,a,t$, но полный rank означает лишь, что
для каждой тройки разрешены три независимых invariant coefficients.

\section{Нормировки не выбирают точку}

Даже если независимо фиксировать weighted sums
\begin{equation}
 \theta_s+\theta_a+3\theta_t=\Theta,
 \qquad
 \kappa_s+\kappa_a+3\kappa_t=K,
 \label{eq:v9-kms-origin-two-normalizations}
\end{equation}
матрица ограничений имеет
\begin{equation}
 \rank A_{\rm norm}=2,
 \qquad
 \dim\ker A_{\rm norm}=4.
 \label{eq:v9-kms-origin-normalized-freedom}
\end{equation}
Таким образом, два trace-like условия оставляют четырёхмерное относительное
семейство. Кроме того, сохраняются две масштабные орбиты
\begin{equation}
 (\beta,\boldsymbol\Delta)\mapsto
 (c\beta,c^{-1}\boldsymbol\Delta),
 \qquad
 (\boldsymbol\kappa,t)\mapsto(c\boldsymbol\kappa,t/c),
 \qquad c>0.
 \label{eq:v9-kms-origin-scale-orbits}
\end{equation}

\section{Два точных нормированных свидетеля}

Рассмотрим два KMS-пакета
\begin{equation}
 \begin{aligned}
 q^{(1)}&=\left(\frac12,\frac13,\frac14\right),
 &\kappa^{(1)}&=(1,1,1),\\
 q^{(2)}&=\left(\frac13,\frac12,\frac14\right),
 &\kappa^{(2)}&=\left(2,1,\frac23\right).
 \end{aligned}
 \label{eq:v9-kms-origin-two-witnesses}
\end{equation}
Они имеют одинаковые weighted normalizations
\begin{equation}
 q_s^{(i)}+q_a^{(i)}+3q_t^{(i)}=\frac{19}{12},
 \qquad
 \kappa_s^{(i)}+\kappa_a^{(i)}+3\kappa_t^{(i)}=5.
 \label{eq:v9-kms-origin-witness-normalizations}
\end{equation}
Оба генератора primitive и имеют Liouvillian rank $35$, но их initial
excitation conductances различны:
\begin{equation}
 R_1=\sum_\alpha d_\alpha\kappa_\alpha^{(1)}q_\alpha^{(1)}
 =\frac{19}{12},
 \qquad
 R_2=\sum_\alpha d_\alpha\kappa_\alpha^{(2)}q_\alpha^{(2)}
 =\frac53,
 \quad d=(1,1,3).
 \label{eq:v9-kms-origin-witness-rates}
\end{equation}
Их faithful стационарное состояниеs также различны:
\begin{equation}
 \rho_1=\frac1{31}\operatorname{diag}(12,6,4,3,3,3),
 \qquad
 \rho_2=\frac1{31}\operatorname{diag}(12,4,6,3,3,3).
 \label{eq:v9-kms-origin-witness-states}
\end{equation}

\section{Почему положительная форма без source недостаточна}

Каноническая multiplicity metric на одной тройке равна
\begin{equation}
 W=\operatorname{diag}(1,1,3),
 \qquad
 \mathbb W=W\oplus W,
 \qquad \det\mathbb W=9,
 \qquad \rank\mathbb W=6.
 \label{eq:v9-kms-origin-common-metric}
\end{equation}
Source-free bounded parent
\begin{equation}
 S_0(\boldsymbol\theta,\boldsymbol\kappa)
 =\frac12\boldsymbol\theta^TW\boldsymbol\theta
 +\frac12\boldsymbol\kappa^TW\boldsymbol\kappa
 \label{eq:v9-kms-origin-source-free-parent}
\end{equation}
имеет единственный минимум
\begin{equation}
 \boldsymbol\theta_*=0,
 \qquad \boldsymbol\kappa_*=0.
 \label{eq:v9-kms-origin-zero-minimum}
\end{equation}
Первое равенство даёт вырожденный Hamiltonian, второе выключает dynamics.
Для положительной внутренней точки необходимы две source-тройки:
\begin{equation}
 S_j=S_0-j_\theta^T\boldsymbol\theta-j_\kappa^T\boldsymbol\kappa,
 \qquad
 \boldsymbol\theta_*=W^{-1}j_\theta,
 \quad
 \boldsymbol\kappa_*=W^{-1}j_\kappa.
 \label{eq:v9-kms-origin-sourced-parent}
\end{equation}
Но задание $(j_\theta,j_\kappa)$ эквивалентно заданию искомых шести чисел,
пока их общий родительский функционал-origin не указан.

\section{Аудит существующих кандидатов}

Проверены семь классов кандидатов:
\begin{equation}
 \begin{aligned}
 \mathcal C_{\rm origin}=\{&\text{symmetry},\text{common trace},
 \text{type observables},\text{KMS},\\
 &\text{Gibbs normalization},\text{phase Laplacian},S_4\}.
 \end{aligned}
 \label{eq:v9-kms-origin-candidate-set}
\end{equation}
Ни один из них не создаёт ненулевые source-тройки без новых коэффициентов:
\begin{equation}
 N_{\rm candidate\ origin}=0/7,
 \qquad
 N_{\rm type\ separation}=3/3,
 \qquad
 N_{\rm normalized\ relative\ freedom}=4.
 \label{eq:v9-kms-origin-candidate-ledger}
\end{equation}
Итоговый физический реестр равен
\begin{equation}
 \begin{aligned}
 N_{\rm gap\ origin}&=0/3,\\
 N_{\rm conductance\ origin}&=0/3,\\
 N_{\rm KMS\ parameter\ origin}&=0/6,\\
 N_{\rm physical\ four\ slot\ parent}&=0/1.
 \end{aligned}
 \label{eq:v9-kms-origin-final-ledger}
\end{equation}

\begin{theorem}[Типовая полнота без численного selector]
Multiplicity, grading и charge образуют невырожденную систему типов для
трёх KMS-каналов. Тем не менее они разрешают весь трёхмерный diagonal
commutant отдельно для gaps и conductances. Две естественные нормировки
оставляют четырёхмерное семейство, содержащее неэквивалентные primitive
generators.
\end{theorem}

\begin{nogocriterion}[Запрет source-free KMS parent]
Положительная invariant quadratic form без линейного source минимизируется
при нулевых gaps и conductances. Ненулевая primitive KMS dynamics требует
двух трёхкомпонентных source-векторов либо механизма, который выводит их из
других slots общего родительского функционала.
\end{nogocriterion}

\begin{protocol}\begin{enumerate}
\item type separation: \textbf{$3/3$};
\item normalized relative freedom: \textbf{$4$};
\item candidate origins: \textbf{$0/7$};
\item gap origin: \textbf{$0/3$};
\item conductance origin: \textbf{$0/3$};
\item KMS parameter origin: \textbf{$0/6$};
\item следующий гейт: \textbf{common source-parent architecture}.
\end{enumerate}\end{protocol}

Машинный сертификат сохранён в
\path{s2t_v9_endpoint_creation_kms_gap_conductance_parent_origin_gate_results.json}.

\begin{thebibliography}{9}
\bibitem{v9-kms-origin-carlen-maas}
E.~A.~Carlen, J.~Maas, \emph{Gradient flow and entropy inequalities for
quantum Markov semigroups with detailed balance}, arXiv:1609.01254.
\bibitem{v9-kms-origin-temme}
K.~Temme, M.~J.~Kastoryano, M.~B.~Ruskai, M.~M.~Wolf, F.~Verstraete,
\emph{The $\chi^2$-divergence and mixing times of quantum Markov processes},
J. Math. Phys. \textbf{51} (2010), 122201.
\end{thebibliography}